\section{Introduction}
\label{sec:intro}
The growing popularity of mobile devices and location-based services has led to an unprecedented influx of geospatial data, capturing detailed information on individual movements, location preferences, and popular routes \citep{Statista2015lbs,cheng2017mobile}. This surge in location data opens up opportunities for a wide range of applications across sectors, from optimizing public transportation systems and enhancing emergency response times to supporting intelligent urban planning and resource allocation \cite{ouchra2023overview}. One key approach in analyzing this data is through frequency analysis, which identifies areas with high concentrations of people over time \cite{geo_hong_21}. By quantifying how often certain locations are visited, frequency analysis provides crucial insights that enable decision-makers to allocate resources effectively and design services that meet population needs.

Frequency analysis, for instance, is widely used for generating heat maps \cite{bao2023mapping,hong2017so,peipei2020study}, which visually represent high-frequency areas or routes based on user interactions. By highlighting patterns in pedestrian traffic, vehicle concentration, or other location-based trends \cite{rikalovic2018spatial}, these maps offer actionable insights into critical areas that could benefit from additional services or infrastructural improvements \cite{miranda2023sok}. Furthermore, tracking frequency changes over time enhances understanding of dynamic trends, allowing analysis of complex, time-dependent phenomena. For instance, observing shifts in traffic flow patterns helps city planners design better road systems and optimize signal timings to alleviate congestion \cite{huang2024festgcn}. In public health, evolving frequency patterns can reveal the spread of infectious diseases \cite{bali2024exploring}, guiding officials to proactively allocate resources where they are most needed. Additionally, analyzing changes in movement patterns informs a range of personalized services and applications, from targeted advertising and location-based recommendations to social behavior research \cite{shi2020study}.

However, while these geospatial analyses offer substantial value across fields, they also raise serious privacy concerns. The data required to generate heat maps over time and track patterns often includes sensitive information about individuals' locations and movements, leading to potential risks related to user identification, behavioral profiling, and location-based discrimination \cite{unnikrishnan2013anonymizing,miranda2023sok,solymosi2023privacy}. As a result, preserving the privacy of users while leveraging the benefits of geospatial data has become a crucial issue. 

Differential Privacy (DP) has emerged as a prominent framework for enabling data analysis while preserving individual privacy \cite{DBLP:conf/icalp/Dwork06,cummings2023advancing}. However, its traditional model relies on a trusted curator, introducing potential vulnerabilities due to the risk of data breaches. To address this limitation, Local Differential Privacy (LDP) has been proposed as a robust alternative. LDP allows individuals to anonymize their data locally before sharing it for analysis, eliminating the need for a centralized curator \cite{DBLP:journals/fttcs/DworkR14,kasiviswanathan2011can}. This decentralized approach has gained significant traction in recent years, with major technology companies such as Apple \cite{DBLP:conf/uss/WangBLJ17}, Google \cite{erlingsson2014rappor}, and Microsoft \cite{DBLP:conf/nips/DingKY17} incorporating LDP into their consumer products.

LDP is particularly well-suited for scenarios involving longitudinal data—data collected from the same individual over time. For instance, in systems where users repeatedly send location reports or other personal information at multiple time intervals, the resulting data forms a longitudinal dataset. LDP ensures privacy protection in such settings by anonymizing each data point before transmission.




Despite the advantages of LDP, current techniques applied to geospatial data often rely on uniform grid structures to partition space for simplicity \cite{andres2013geo,ye2021collecting}, which can lead to an inadequate representation of the data. These approaches apply the same level of granularity to all regions, regardless of the density of points of interest, potentially compromising both the accuracy of the analyses and the utility of the aggregated data. This issue, known as non-uniformity error, arises when the assumption of uniform data distribution within grid cells does not align with the actual data distribution \cite{qardaji2013differentially}. In dense regions, this mismatch can lead to underrepresentation of critical areas, introducing bias into the results, while in sparse regions, excessive granularity can amplify noise, reducing the utility of the data. 
Because these methods don't consider how data is distributed within a region, they can produce large errors.
For example, in dense regions, uniform grids may be too coarse, leading to significant bias when answering queries. On the other hand, regions with sparse data may be mapped by a fine-grained grid, resulting in weak utility and high noise during query answering \cite{qardaji2013differentially}. \looseness=-1

In this paper, we propose a new approach based on adaptive, non-uniform grids that dynamically adjust to data point densities, allowing for greater accuracy in regions with high individual concentration and coarser granularity in less dense areas. Our methodology combines the flexibility of adaptive grids with the privacy guarantees offered by Local Differential Privacy. By iteratively refining cell sizes as new data is reported, our approach incrementally adjusts to temporal variations in data distribution. This temporal refinement is essential to accurately capture the correlations present among consecutive reports submitted under Local Differential Privacy (LDP), enabling our framework to deliver more precise and context-aware analyses. Thus, these enhancements make our framework, called ALOG (Adaptive Longitudinal Grids), particularly well-suited for addressing the challenges of privacy-preserving geospatial data analysis.
% 
\subsubsection*{Contributions} In summary, this paper makes the following contributions:\looseness=-1
\begin{itemize}
    \item We propose ALOG, an adaptive grid-based framework for LDP-compliant longitudinal data collection, which exploits temporal correlations and dynamically refines spatial partitions to balance privacy and utility. 
    \item  We analyze various grid refinement strategies to detect and adapt to changing patterns in data density throughout the grid adaptation windows with a detailed investigation of the factors that influence the utility-privacy trade-off.
    \item To assess the effectiveness of our adaptive longitudinal data collection framework, we conduct comprehensive evaluations using both synthetic and real-world datasets while comparing them with different prior state-of-the-art techniques.\looseness=-1
\end{itemize}

\subsubsection*{Organization} Section 2 formalizes the problem tackled in this work, and presents the foundations and essential definitions used throughout the work. In Section 3, we present
the related work and discuss how the main existing approaches
work and their limitations. We present our approach in Section 4. Experimental results are shown in Section 5. Finally, Section 6 concludes
this work.
